% !TeX root = ../深度学习与大语言模型技术原理.tex
% 由主文档通过 \input 载入；请编译主文档。
\setcounter{chapter}{8} % 保留原书章号。
\chapter{分布式训练与集合通信}
\label{ch:distributed}

\section{集合通信及其成本}

集合通信（collective communication）让一组设备协调交换张量。AllReduce 对各设备张量做归约，并使各设备得到完整结果；AllGather 把分片收集成完整张量；ReduceScatter 先归约再把结果分片；All-to-All 让各设备向其他设备发送不同数据块。点对点（point-to-point）通信用于指定设备之间的发送和接收。NCCL 提供 GPU 集合通信实现，算法效果依赖互连、拓扑和消息大小\cite{nccl}。

一个粗略通信模型为 $T_{\mathrm{comm}}\approx a\tau+S/\mathcal B$，其中 $a$ 为通信阶段数，$\tau$ 为每阶段启动延迟，$S$ 为关键路径传输字节数，$\mathcal B$ 为有效带宽。小消息容易受延迟限制，大消息容易受带宽限制。把梯度分桶可以减少小消息开销，并与反向计算重叠；桶太大又会推迟通信启动。只看峰值网络带宽不足以预测训练速度。

\section{数据并行与正确的梯度平均}

数据并行（data parallelism, DP）在每张设备复制模型，各设备处理不同数据，再同步梯度。DistributedDataParallel（DDP）通常在反向过程中同步梯度，各设备执行相同优化器更新以保持参数一致\cite{ddp}。模型副本相同不意味着数据应相同；重复采样同一批会浪费并行能力，并改变有效批量理解。

若有 $D$ 张设备，每卡损失和为 $S_r$、有效词元数为 $N_r$，且通信对梯度取平均，则各卡可反向传播 $D S_r/N_{\mathrm{all}}$，其中 $N_{\mathrm{all}}=\sum_rN_r$。平均后的梯度为
\[
\frac1D\sum_r\nabla_\theta\frac{D S_r}{N_{\mathrm{all}}}
=\nabla_\theta\frac{\sum_rS_r}{N_{\mathrm{all}}}.
\]
这说明软件的归约约定必须与损失归一化配套。若每卡先求局部平均而词元数不同，则得到的是设备等权目标。

DDP 不直接减少每卡的参数和优化器状态。多卡能扩大吞吐与有效批量，但模型单卡装不下时，需要分片或模型并行。检查恢复时，还应保持采样器状态与全局步数一致，避免重复或遗漏数据。

\section{ZeRO 与完全分片数据并行}

\emph{零冗余优化器（Zero Redundancy Optimizer, ZeRO）}按阶段分片状态\cite{zero}。ZeRO-1 分片优化器状态；ZeRO-2 进一步分片梯度；ZeRO-3 再分片参数。完全分片数据并行（Fully Sharded Data Parallel, FSDP）采用相关思想，在需要计算某个参数块时聚合参数，反向后归约并重新分片，具体策略依实现而定。

设完整参数、梯度和优化器状态分别占 $M_p,M_g,M_o$，忽略临时聚合时，每卡静态状态近似如下：
\begin{center}
\begin{tabular}{ll}
\toprule
方案 & 每卡静态状态 \\
\midrule
普通数据并行 & $M_p+M_g+M_o$ \\
ZeRO-1 & $M_p+M_g+M_o/D$ \\
ZeRO-2 & $M_p+(M_g+M_o)/D$ \\
ZeRO-3 & $(M_p+M_g+M_o)/D$ \\
\bottomrule
\end{tabular}
\end{center}
表中的估算没有包含激活和当前聚合参数块，因此不是峰值显存。更细粒度分片降低聚合峰值，却可能增加通信次数；CPU 或存储卸载减少设备占用，同时引入传输与调度成本。

\section{张量并行与流水线并行}

张量并行（tensor parallelism, TP）把单层矩阵运算分到多卡。以行批量约定 $Y=XW$ 为例，按 $W$ 的输出列分块 $W=[W_1,\ldots,W_D]$，各卡计算 $Y_r=XW_r$，拼接得到 $Y$；这是列并行。若按输入维分块，$X=[X_1,\ldots,X_D]$ 且 $W$ 沿行分块，则 $Y=\sum_rX_rW_r$，需归约各卡部分结果；这是行并行。实际 Transformer 常配对使用，避免不必要的中间聚合。

流水线并行（pipeline parallelism, PP）把不同层放在不同阶段，微批量依次流过。以 $p$ 个等耗时阶段、$m$ 个微批量的简单前向流水线为例，执行槽数为 $m+p-1$，相对理想满载的利用率为 $m/(m+p-1)$。完整训练的气泡还依赖前反向调度；一前向一反向（1F1B）能改善激活占用和调度，但不能据简单公式断言所有情形的利用率。

交错流水线把每卡的层拆成多个虚拟阶段，以减少某些调度气泡，同时增加通信与实现复杂度。TP 倾向需要频繁低延迟通信，PP 通常传输阶段边界激活；如何跨节点部署，应根据拓扑和层负载决定。模型能放下以后，再增加切分未必能加速。

\section{上下文并行、专家并行与训练框架}

上下文并行（context parallelism, CP）沿序列维度拆分词元，使每卡保存和计算部分位置的激活。对于全注意力，本地查询仍需要访问相关的全局键值，因此必须交换键值块或采用环式调度；仅把序列截断给不同卡而不通信，会改变模型可见上下文。因果掩码带来位置间负载差异，需要有意识地分配序列块。

专家并行按第~\ref{ch:moe} 章的路由分配专家，与 CP 的序列划分、TP 的层内划分和 PP 的层间划分不同。多维并行会形成不同通信组，各组的设备数乘积受总设备数约束；并非所有维度都必须开启。先由显存约束确定最小必要切分，再优化吞吐，通常更容易解释性能。

DeepSpeed 集成 ZeRO 等训练机制，Megatron-Core 组织大模型层与多维并行，PyTorch 提供基础张量与分布式能力，Accelerate 辅助组织设备和训练流程。后训练框架如 verl、OpenRLHF 还需协调生成、奖励、价值估计与策略更新。它们不会自动消除错误的损失掩码、过时的采样策略或不正确的跨卡归一化；框架迁移应比较同一小任务上的目标、梯度与有效词元统计。

